\documentclass[11pt,a4paper]{article}
\usepackage[margin=22mm]{geometry}
\usepackage{amsmath,amssymb,booktabs,microtype}
\usepackage[hidelinks]{hyperref}
\newcommand{\ii}{\mathrm i}
\newcommand{\ket}[1]{|#1\rangle}
\newcommand{\bra}[1]{\langle#1|}
\newcommand{\tr}{\operatorname{tr}}
\newcommand{\Var}{\operatorname{Var}}
\title{Route G-R14: One Event Rule for Clock Rate,\\Loop Phase and Coherence}
\author{A finite candidate with a common-rate test and an explicit memory rule}
\date{15 October 2026}
\begin{document}
\maketitle
\begin{abstract}
One binary event register couples to the sum of all probe charges.
A single occupation probability determines clock drift, phase noise,
and a three-link interference amplitude. They cannot be fitted independently.
The common mean event rate is not generally a common fractional clock rate;
the nontrivial exact common-rate case is balanced occupation, within a
readable phase band. A rank-two Gram identity locks loop phase to visibility.
Keeping the same register around the loop instead of using fresh registers
changes the answer from a nontrivial complex amplitude to exactly one.
This identifies inter-event memory as physical information missing from a
single scalar density. It does not derive physical time or the original
paper's full theory.
\end{abstract}

\section{Inputs and one microscopic rule}
Chen's event accumulation and compatible comparisons motivate this
construction~\cite{chen}; its Hilbert spaces, probabilities and dynamics
are additional postulates. Three links label comparisons, not spatial
dimensions or rotations of matter. We retain the G-R9--G-R13 results
unchanged; this is a new bounded candidate, not a new control pulse.

Each ordered event has a two-state register $R$ and a probe $S$:
\begin{equation}
 \ket{r_p}=\sqrt{1-p}\ket0+\sqrt p\ket1,\quad
 P=\ket1\bra1,\quad
 U(Q)=\exp(-\ii a Q\otimes P),\quad a=2\pi/3.\label{rule}
\end{equation}
The primitive three-cycle phase $a$ is chosen, not measured or derived.
$p\in[0,1]$ is the only varying response parameter. $Q$ is a specified
dimensionless probe generator: $Q=q\ket1\bra1$ for a clock, or
$Q=\ket{\mathrm{loop}}\bra{\mathrm{loop}}$ for an interferometer.
For simultaneous probes use their \emph{sum}, with identities on other
factors. Coupling all probes this way is an explicit universality
postulate; it is not a consequence of naming the register an event.
Fix two clock charges $q_A=1/4,q_B=1/2$, not particle masses.

Fresh registers start in a product state and are retained but unread.
Tracing one gives the exact completely positive trace-preserving map
\begin{align}
 \mathcal E_p(\rho)&=(1-p)\rho+p e^{-\ii aQ}\rho e^{\ii aQ},\label{channel}\\
 F_p(x)&=1-p+p e^{-\ii x},&
 (\mathcal E_p^n\rho)_{jk}&=F_p\big(a(q_j-q_k)\big)^n\rho_{jk}.\label{F}
\end{align}
Indeed $U=I\otimes\ket0\bra0+e^{-\ii aQ}\otimes\ket1\bra1$;
its two Kraus operators have $\sum K_\ell^\dagger K_\ell=I$.
No loss branch is discarded. Ordered $n$, register supply, quantum
probabilities and the measurement reference remain inputs, not seconds
or a derived arrow of time. This is a repeated-interaction model of a
standard kind~\cite{collision}; no new general quantum formalism is claimed.

\newpage
\section{Mean event progress is not automatically clock frequency}
Reading all $P_\ell$ would give $K_n=\sum_{\ell=1}^n P_\ell$ with
\begin{equation}
 \mathbb E K_n=np,\qquad \Var K_n=np(1-p),\qquad
 N:=\mathbb E K_n/n=p.\label{mean}
\end{equation}
The probe interaction preserves these register probabilities.
The same $K_n$ acts on every probe present at an event. Thus $N$ is a
genuine common \emph{mean active-event fraction}, not a frequency defined
from a trace or arbitrarily inserted into three unrelated formulas.

Start a clock in $(\ket0+\ket1)/\sqrt2$. Its two quadratures measure
\begin{equation}
 C_q(n):=2\rho_{10}(n)=F_p(aq)^n,
 \qquad V_q(n)=|F_p(aq)|^n.\label{clock}
\end{equation}
On the continuous phase branch connected to $x=aq=0$, define
\begin{equation}
 \vartheta_q=-\arg F_p(aq),\quad
 R_q=\frac{\vartheta_q}{aq},\quad
 d_q=-\log|F_p(aq)|.\label{rates}
\end{equation}
$\vartheta_q$ is phase advance per update, $R_q$ its fraction of the
bare advance $aq$, and $d_q$ visibility loss per update. These are not
rates per laboratory second. If $F_p=0$, phase is undefined. From
Eq.~\eqref{F},
\begin{align}
 |F_p(x)|^2&=1-4p(1-p)\sin^2(x/2),\label{modulus}\\
 \log F_p(x)&=-\ii px-\tfrac12p(1-p)x^2
 +\tfrac{\ii}{6}p(1-p)(1-2p)x^3+O(x^4),\\
 R_q&=p-\frac{p(1-p)(1-2p)}6(aq)^2+O((aq)^4),\label{rateexp}\\
 d_q&=\tfrac12p(1-p)(aq)^2+O((aq)^4).
\end{align}
The cumulants follow by differentiating $\log(1-p+pe^t)$ at $t=0$;
substitute $t=-\ii x$. Reality gives odd phase and even log modulus,
which fixes the remainder orders. The mean, phase drift and noise all
come from the same Bernoulli distribution, without an extra damping law.

\paragraph{Exact common-rate proposition.}
If $R_q$ is constant for all $q$ in an open neighborhood of zero, its
constant must be $p$ by the limit in Eq.~\eqref{rateexp}. Its quadratic
coefficient must vanish, so $p=0,1/2,1$. Conversely $p=0,1$ are
deterministic, while
\begin{equation}
 F_{1/2}(x)=e^{-\ii x/2}\cos(x/2),\qquad
 R_q=\tfrac12\quad (|aq|<\pi).\label{balanced}
\end{equation}
This proves necessity and sufficiency on a connected readable band.
The $p=0$ clock does not advance. At $p=1/2$, visibility remains
$|\cos(aq/2)|^n$, not one; zeros and subsequent phase jumps are not
time reversal. Equal occupation is not selected by dynamics here.
For generic $p$, common mean progress yields only approximate low-charge
universality. With resolved phase winding and $\vartheta_A\ne0$, one
can eliminate the unobserved count:
$\phi_B=(\vartheta_B/\vartheta_A)\phi_A$ and
$\log V_B=(d_B/d_A)\log V_A$ when $d_A\ne0$.

\newpage
\section{The same rule fixes a closed-loop phase and its visibility}
Set $V=e^{-\ii aP}$, so $V^3=I$. Define three record rays
\begin{equation}
 \ket{e_j}=V^j\ket{r_p},\quad j=0,1,2,\qquad
 z=\langle e_0|e_1\rangle=\langle e_1|e_2\rangle
   =\langle e_2|e_0\rangle=F_p(a).
\end{equation}
Writing $z=r e^{\ii\theta}$, the oriented product is
\begin{equation}
 W=z^3=r^3e^{\ii\Phi},\qquad \Phi=3\theta\pmod{2\pi}.
 \label{loop}
\end{equation}
Rephase each ray by $e^{\ii\chi_j}$. Each overlap gains
$e^{\ii(\chi_{j+1}-\chi_j)}$ and their closed product is unchanged.
Reversing orientation conjugates $W$. This is a three-state Bargmann
product~\cite{bargmann}, not by itself a magnetic field.

The Gram matrix is
\begin{equation}
 G=\begin{pmatrix}1&z&z^*\\z^*&1&z\\z&z^*&1\end{pmatrix}\succeq0.
\end{equation}
There are three vectors in a two-dimensional space, hence
\begin{equation}
 \boxed{\det G=1-3r^2+2r^3\cos\Phi=0},\qquad
 r^2=1-3p(1-p).\label{constraint}
\end{equation}
The determinant follows directly by expanding its six permutation
terms. The second identity follows from
$\cos a=-1/2$ in Eq.~\eqref{modulus}. Thus $1/2\le r\le1$;
unit link visibility forces $W=1$ in this family. In higher record
dimension the general Gram positivity condition is $\det G\ge0$,
not the equality above. The equality is therefore a falsifiable
rank-two model restriction, not a universal law of phases.

\paragraph{Operational loop protocol with all outputs.}
Prepare a path qubit $(\ket{\mathrm{ref}}+\ket{\mathrm{loop}})/\sqrt2$
and three fresh registers. On each oriented link apply Eq.~\eqref{rule}
with $Q=\ket{\mathrm{loop}}\bra{\mathrm{loop}}$. The reference sees
identity; the loop sees $V$ on each successive register. Its final state is
\begin{equation}
 \ket\Psi=\frac{1}{\sqrt2}\left(
 \ket{\mathrm{ref}}\ket{r_p}^{\otimes3}
 +\ket{\mathrm{loop}}(V\ket{r_p})^{\otimes3}\right).
\end{equation}
Tracing registers gives $2\rho_{\mathrm{loop,ref}}=W$ and equal
path populations $1/2$. At analysis phase $\varphi$ the two output
probabilities are
\begin{equation}
 P_\pm(\varphi)=\tfrac12\left[1\pm
       \Re(e^{-\ii\varphi}W)\right],\quad
 P_++P_-=1,\quad \mathcal V_{\rm loop}=r^3.\label{fringe}
\end{equation}
This realizes Eq.~\eqref{loop} without filtering, heralding or assigning
a separate coherence factor. Spatial transport and beam-splitter
hardware are not derived by naming the controlled comparisons links.

\newpage
\section{Memory composition and a joint prediction}
One-event data do not fix multi-event data without a memory rule.
Reuse a single register on all three links, without reading or resetting
it, while keeping the identical gate and initial register state. Then
\begin{equation}
 W_{\rm reused}=\langle r_p|V^3|r_p\rangle=F_p(3a)=1,
 \qquad W_{\rm fresh}=F_p(a)^3.\label{memory}
\end{equation}
For $p=1/2$, fresh records give $(\Phi,\mathcal V)=(\pi,1/8)$;
reuse gives $(0,1)$. The elementary marginals and mean progress $N$
are identical. For a general sequence the relevant object is
\begin{equation}
 \mathcal F(x_1,\ldots,x_n)=
 \left\langle e^{-\ii\sum_\ell x_\ell P_\ell}\right\rangle.
\end{equation}
Independent registers factorize this into $\prod F_p(x_\ell)$;
the reused register gives $F_p(\sum x_\ell)$. Replacing the latter
by the former silently changes the physics. In a classical diagonal
description this is exactly the distinction between independent and
perfectly correlated binary variables. Nontrivial phase of an averaged
loop is not proof that each microscopic branch has the same holonomy.

\paragraph{Two clocks share noise, not independent damping constants.}
Take $Q=q_A P_A+q_B P_B$ on two clocks. In the basis
$00,01,10,11$, its charges are $0,q_B,q_A,q_A+q_B$. From an initial
product plus state, Eq.~\eqref{F} gives
\begin{equation}
 C_A=F_p(aq_A)^n,\quad C_B=F_p(aq_B)^n,\quad
 4\rho_{10,01}=F_p(a(q_A-q_B))^n.\label{relative}
\end{equation}
For $q_A=q_B$ the relative channel has unit normalized coherence even
when each local dial fades. A comparison of identical clocks alone
can therefore miss a common fluctuation. The state of two clocks is a
mixture of product phase shifts, so this example does not create
clock--clock entanglement from the product input.

\paragraph{Clock and loop measured together.}
Now use $Q=qP_C+P_{\rm loop}$ at each of three events. The four-charge
list is $0,1,q,q+1$. The joint double-off-diagonal element is
\begin{equation}
 J:=4\rho_{11,00}=F_p(a(q+1))^3
 \ne C_q(3)W_{\rm fresh}\quad\hbox{in general}.\label{joint}
\end{equation}
It is measurable from joint quadratures, since
\begin{equation}
 \langle XX-YY\rangle=\Re J,\qquad
 \langle YX+XY\rangle=\Im J.
\end{equation}
Both identities follow by expanding the Pauli matrices; the
$\rho_{10,01}$ terms cancel. At $p=1/4,q=1/4$,
$J=0.126833033-0.104777678\ii$, whereas the incorrect product of
marginals is $0.049144099-0.263300461\ii$.
This discrepancy is not an adjustable loop or noise parameter.
All probes in an event use one shared register; separate experiments
use identically prepared ones.

\paragraph{What loss of visibility means.}
The global state remains pure and the inverse global unitary restores
the initial state. Tracing unread registers is loss of access, not a
proof of irreversible decoherence. Reset, isolation, a persistent record
and controller work remain implementation questions. The gate commutes
with bare $Q$ and with every $P_\ell$; its phase effects alone are not
evidence of an energy transfer or changed particle rest mass.

\newpage
\section{Numerical checks and the research boundary}
The verifier constructs full system--register unitaries (up to
$4\times2^3=32$ states), takes partial traces, and compares them
with iterated Kraus maps and the analytic characteristic function.
It checks all probabilities, complete positivity, register count
moments, conserved bare charge, loop rephasing, rank-two Gram identity,
memory reuse, joint quadratures, relative-clock protection and the
small-argument expansion. The fixtures below are chosen rational
occupations, not a parameter search or a fit to observed time dilation.
All 223 new checks and 6672 checks across 27 verifiers pass;
G-R9/G-R13 code, cards and results keep their hashes.

\begin{center}
\begin{tabular}{rrrrr}\toprule
$p=N$&$R_A$&$R_B$&$\Phi$ (rad)&$\mathcal V_{\rm loop}$\\\midrule
$1/4$&0.245641993&0.231631437&$-1.000419517$&0.289379050\\
$1/2$&0.500000000&0.500000000&$\pi$&0.125000000\\
$3/4$&0.754358007&0.768368563&$+1.000419517$&0.289379050\\\bottomrule
\end{tabular}
\end{center}
At $p=1/2$ the two clock visibilities after three events are
0.901221065 and 0.649519053 despite their identical fractional rate.
For generic $p$, identifying $N=p$ with both measured $R_A,R_B$ is
rejected, not repaired by fitting the charges. The endpoint fixtures
$p=0,1$ are deterministic controls. Numerical tolerances test the exact
identities; they are not the scientific result.

\paragraph{What has been completed.}
A single specified rule and register state give three distinct,
jointly constrained outputs: mean progress and actual clock phase,
closed-loop phase, and coherence. A nontrivial exact common-rate
case exists, but its occupation balance and finite phase band are
essential. The proposed arbitrary-$p$ universal clock interpretation
fails inside the same model. The model supplies a stronger statement
than three independently chosen response functions, not an unconditional
proof of Chen's claims.

\paragraph{What remains an input.}
No law selects $p$, $a$, event ordering or the sum-of-charges coupling.
No physical $E_i/E_d$, seconds, particle spectrum, Maxwell field,
gravity, or new Bell bound follows. A relation $p(E_d)$ inserted later
would be a new hypothesis, not a result of these calculations. The
reduced commuting-probe predictions also admit a classical correlated
binary-noise realization. Distinguishing that from a genuinely quantum
event mechanism requires additional noncommuting record probes.

\paragraph{Next bounded choices.}
First constrain the memory rule with a three-link comparison and the
joint observable in Eq.~\eqref{joint}, retaining the same $p$ inferred
from register occupations. A failure of the Gram equality, the joint
response, or the predicted rate pair falsifies this candidate with its
stated preparation, not quantum theory. Alternatively, if exact common
fractional time across a varying environment is indispensable, require
a physical principle selecting or protecting the balanced rule; do not
silently equate mean rate and phase rate. Controller calibration from
G-R13 remains a side branch and does not block either question. No new
pulse scan, Route-F promotion or long prerequisite proof is introduced.

\begingroup\small
\begin{thebibliography}{9}
\bibitem{chen} B. Chen, \emph{Energy spectrum for elementary particles},
final version (2022), \href{https://indico.cern.ch/event/1109513/contributions/4969967/attachments/2480973/4410560/Energy\%20spectrum\%20for\%20elementary\%20particles_final.pdf}{CERN Indico source}.
Conceptual inspiration only; source receipt and prior section audit are in
\texttt{route\_g/SOURCES.md}.
\bibitem{collision} F. Ciccarello, S. Lorenzo, V. Giovannetti and G. M. Palma,
\emph{Quantum collision models: Open system dynamics from repeated interactions},
Phys. Rep. 954, 1 (2022), \href{https://arxiv.org/abs/2106.11974}{arXiv:2106.11974}.
\bibitem{bargmann} E. M. Rabei, Arvind, R. Simon and N. Mukunda,
\emph{Bargmann invariants and geometric phases: A generalized connection},
Phys. Rev. A 60, 3397 (1999),
\href{https://arxiv.org/abs/quant-ph/9907080}{arXiv:quant-ph/9907080}.
\end{thebibliography}
\endgroup
\end{document}
